% conference Conference Paper - System Architecture
% BanditGPT: Prior Initialization and Latent Semantic Transfer

\documentclass[sigconf]{acmart}

\usepackage{amsmath}
\usepackage{amssymb}
\usepackage{algorithm}
\usepackage{algorithmic}
\usepackage{booktabs}
\usepackage{graphicx}

% Remove copyright box for technical documentation
\settopmatter{printacmref=false}
\renewcommand\footnotetextcopyrightpermission[1]{}
\pagestyle{plain}

\begin{document}

\title{BanditGPT System Architecture: \\
Prior Initialization and Latent Semantic Transfer}

\author{BanditGPT Research Team}
\affiliation{%
  \institution{Production LLM Router System}
}

\begin{abstract}
We present a production-grade contextual bandit router that eliminates cold-start inefficiency through multi-tiered prior initialization and latent semantic transfer. Our approach achieves 92\% cost reduction at production quality by transferring learned preferences ($\boldsymbol{\theta}$) rather than raw covariance structures ($\mathbf{A}$), avoiding the ``confident transfer trap'' that causes premature convergence in traditional Bayesian transfer methods. Validated on 1,871 real prompts with a domain shift of PSI=0.42, our system demonstrates robust recovery through T-shirt sizing injection and semantic DNA matching.
\end{abstract}

\maketitle

\section{Introduction}

Production deployment of contextual bandits for LLM routing presents unique challenges: (1) cold-start inefficiency when new models are introduced, (2) negative transfer from domain mismatch between training and production distributions, and (3) the need to incorporate business logic (e.g., cost constraints) into learned policies. This paper presents the architectural innovations in BanditGPT that address these challenges through a mathematically grounded multi-tiered initialization strategy.

\section{Background}

\subsection{Contextual Bandits for LLM Routing}

We model LLM routing as a contextual bandit problem where:
\begin{itemize}
    \item \textbf{Context} $\mathbf{x}_t \in \mathbb{R}^d$: Semantic embedding of prompt $t$
    \item \textbf{Arms} $\mathcal{M} = \{m_1, \ldots, m_K\}$: Set of $K$ available LLMs
    \item \textbf{Reward} $r_t \in [0, 1]$: Quality score for selected model
\end{itemize}

The router maintains per-arm parameters $(\mathbf{A}_m, \mathbf{b}_m)$ for each model $m \in \mathcal{M}$ and selects arms using LinUCB~\cite{li2010contextual}:

\begin{equation}
m_t = \argmax_{m \in \mathcal{M}} \left[ \boldsymbol{\theta}_m^\top \mathbf{x}_t + \alpha \sqrt{\mathbf{x}_t^\top \mathbf{A}_m^{-1} \mathbf{x}_t} \right]
\label{eq:linucb}
\end{equation}

where $\boldsymbol{\theta}_m = \mathbf{A}_m^{-1} \mathbf{b}_m$ represents learned preferences.

\subsection{The Cold-Start Problem}

Standard LinUCB initialization uses identity matrices ($\mathbf{A}_m = \lambda \mathbf{I}$, $\mathbf{b}_m = \mathbf{0}$), resulting in:
\begin{itemize}
    \item High initial uncertainty ($\mathbf{x}^\top \mathbf{A}_m^{-1} \mathbf{x} \approx d/\lambda$)
    \item Random exploration during burn-in phase
    \item 10$\times$ sample inefficiency compared to warm-start
\end{itemize}

For production LLM routing, this translates to thousands of suboptimal routing decisions before convergence.

\section{Prior Initialization and Latent Semantic Transfer}

The robustness of banditGPT in production environments relies on a \textbf{multi-tiered initialization strategy} that eliminates the ``Cold Start'' problem while protecting against ``Negative Transfer''. This section details the mathematical justification for transferring learned preferences ($\boldsymbol{\theta}$) rather than raw covariance structures ($\mathbf{A}$).

\subsection{The ``Confident Transfer Trap''}

Standard Bayesian transfer methods often inherit the full posterior of a source model. If a source model (e.g., a legacy flagship) has served millions of requests, its precision matrix $\mathbf{A}$ will have high eigenvalues, resulting in extremely narrow confidence intervals:

\begin{equation}
\mathbf{A}_{\text{source}} = \sum_{i=1}^{N_{\text{large}}} \mathbf{x}_i \mathbf{x}_i^\top + \lambda \mathbf{I}
\end{equation}

where $N_{\text{large}} \gg 10^6$ requests.

\paragraph{Naive Transfer (Problematic):}
\begin{align}
\mathbf{A}_{\text{new}} &= \alpha \mathbf{A}_{\text{source}} \label{eq:naive_transfer_A} \\
\mathbf{b}_{\text{new}} &= \alpha \mathbf{b}_{\text{source}} \label{eq:naive_transfer_b}
\end{align}

where $\alpha \in (0, 1]$ is a dampening factor.

\paragraph{Problem:} Even with $\alpha = 0.1$, the new model inherits:
\begin{equation}
\text{Effective samples} = \alpha N_{\text{large}} \approx 100{,}000
\end{equation}

This causes the router to \textbf{prematurely conclude} the new model is sub-optimal before it has sufficient data to prove its unique utility. The confidence intervals remain tight:

\begin{equation}
\text{Var}[\boldsymbol{\theta}_{\text{new}}^\top \mathbf{x}] = \mathbf{x}^\top \mathbf{A}_{\text{new}}^{-1} \mathbf{x} \approx \frac{\mathbf{x}^\top \mathbf{x}}{\alpha N_{\text{large}}} \ll \frac{\mathbf{x}^\top \mathbf{x}}{\lambda}
\end{equation}

resulting in insufficient exploration and ``fossilization'' of the routing policy.

\subsection{$\boldsymbol{\theta}$-Transfer via Identity Reset}

To solve this, we implement a \textbf{Latent Semantic Transfer} protocol within \texttt{router.py} that separates preference transfer (what the model is good at) from confidence transfer (how certain we are).

\subsubsection{Semantic DNA Matching}

When a new model is registered, we embed its metadata (speed, capabilities, and ID) into a semantic ``DNA'' string:

\begin{equation}
\text{DNA}(m) = \texttt{encode}\left( \text{id}(m) \oplus \text{capabilities}(m) \oplus \text{speed}(m) \right)
\end{equation}

where $\oplus$ denotes string concatenation. We then identify the nearest neighbor in the existing registry using cosine similarity:

\begin{equation}
m^* = \argmax_{k \in \mathcal{M}_{\text{existing}}} \frac{\text{DNA}(m_{\text{new}})^\top \text{DNA}(m_k)}{\|\text{DNA}(m_{\text{new}})\| \|\text{DNA}(m_k)\|}
\label{eq:neighbor_selection}
\end{equation}

\paragraph{Example:} For a new model ``anthropic/claude-3.5-haiku'' with capabilities [``coding'', ``general''] and speed ``fast'', the DNA string becomes:
\begin{verbatim}
"anthropic claude 3.5 haiku coding general fast"
\end{verbatim}

This is embedded using a sentence transformer~\cite{reimers2019sentence} and matched to existing models. Typically, ``anthropic/claude-3-sonnet'' emerges as the nearest neighbor with similarity $> 0.85$.

\subsubsection{Preference Extraction}

We extract the neighbor's learned preference vector, which encapsulates the model's performance manifold across the semantic space:

\begin{equation}
\boldsymbol{\theta}_{\text{neighbor}} = \mathbf{A}_{m^*}^{-1} \mathbf{b}_{m^*}
\label{eq:preference_extraction}
\end{equation}

This vector $\boldsymbol{\theta}_{\text{neighbor}} \in \mathbb{R}^d$ represents the neighbor's learned utility function over the $d$-dimensional semantic space (typically $d=24$ after PCA compression).

\subsubsection{A-Matrix Reset}

\textbf{Critical Innovation:} We initialize the new model's $\mathbf{A}$ matrix as a \textbf{fresh identity matrix} scaled by the regularization parameter:

\begin{equation}
\mathbf{A}_{\text{new}} = \lambda \mathbf{I}
\label{eq:A_reset}
\end{equation}

This ensures the model begins with \textbf{maximum uncertainty} and \textbf{high exploration potential}. The confidence intervals are wide:

\begin{equation}
\text{Var}[\boldsymbol{\theta}_{\text{new}}^\top \mathbf{x}] = \mathbf{x}^\top (\lambda \mathbf{I})^{-1} \mathbf{x} = \frac{\|\mathbf{x}\|^2}{\lambda}
\end{equation}

allowing the new model to quickly diverge from the neighbor if its performance differs.

\subsubsection{Prior Strength Scaling}

The new $\mathbf{b}$ vector is initialized as:

\begin{equation}
\mathbf{b}_{\text{new}} = \lambda \cdot \boldsymbol{\theta}_{\text{neighbor}} \cdot n_{\text{effective}}
\label{eq:b_initialization}
\end{equation}

where $n_{\text{effective}} \in \mathbb{R}^+$ simulates a pseudocount of observations to provide a stable initial ``anchor'' for the model's performance without fossilizing its policy.

\paragraph{Dynamic $n_{\text{effective}}$ Scaling:}
The prior strength is adjusted based on semantic similarity:

\begin{equation}
n_{\text{effective}} = \begin{cases}
5.0 & \text{if } \text{sim}(m_{\text{new}}, m^*) > 0.8 \quad \text{(strong transfer)} \\
3.0 & \text{if } 0.6 < \text{sim}(m_{\text{new}}, m^*) \leq 0.8 \quad \text{(moderate)} \\
1.0 & \text{if } \text{sim}(m_{\text{new}}, m^*) \leq 0.6 \quad \text{(weak transfer)}
\end{cases}
\label{eq:n_effective}
\end{equation}

This ensures:
\begin{itemize}
    \item \textbf{High similarity} ($> 0.8$): Trust neighbor's preferences strongly
    \item \textbf{Low similarity} ($< 0.6$): Weak prior, maximum exploration
\end{itemize}

\subsubsection{Mathematical Properties}

The $\boldsymbol{\theta}$-only transfer preserves preference direction while resetting confidence:

\begin{align}
\boldsymbol{\theta}_{\text{new}} &= \mathbf{A}_{\text{new}}^{-1} \mathbf{b}_{\text{new}} \\
&= (\lambda \mathbf{I})^{-1} (\lambda \cdot \boldsymbol{\theta}_{\text{neighbor}} \cdot n_{\text{effective}}) \\
&= \frac{1}{\lambda} \cdot \lambda \cdot \boldsymbol{\theta}_{\text{neighbor}} \cdot n_{\text{effective}} \\
&= n_{\text{effective}} \cdot \boldsymbol{\theta}_{\text{neighbor}}
\label{eq:theta_preservation}
\end{align}

Thus, the new model's initial predictions are:
\begin{equation}
\hat{r}_{\text{new}}(\mathbf{x}) = n_{\text{effective}} \cdot \boldsymbol{\theta}_{\text{neighbor}}^\top \mathbf{x}
\end{equation}

scaled by $n_{\text{effective}}$ to control prior strength.

\subsection{Business Logic Injection}

Finally, we apply \textbf{T-Shirt Sizing} biases (e.g., \texttt{slow\_bias}, \texttt{fast\_bias}) from the \texttt{RouterConfig}. These biases are scaled by the current confidence (the diagonal of $\mathbf{A}$) to ensure that human-intuition---such as ``expensive models are for reasoning''---is represented with sufficient mathematical weight to guide the early exploration of the 5.8\% Hard Cluster versus the 94.2\% Easy Cluster.

\subsubsection{Confidence-Scaled Bias Injection}

After loading warmup priors or performing semantic transfer, we adjust the bias term:

\begin{equation}
b_m[\text{bias}] \leftarrow b_m[\text{bias}] + A_m[\text{bias}, \text{bias}] \times \Delta_{\text{speed}}
\label{eq:bias_injection}
\end{equation}

where $\Delta_{\text{speed}}$ depends on the model's speed profile:

\begin{equation}
\Delta_{\text{speed}} = \begin{cases}
+0.5 & \text{if } \text{speed}(m) = \text{``fast''} \\
-0.5 & \text{if } \text{speed}(m) = \text{``slow''} \\
0.0 & \text{otherwise}
\end{cases}
\label{eq:speed_bias}
\end{equation}

\paragraph{Why Confidence Scaling?}
After warmup with 80k battles, $b[\text{bias}] \approx 1000$ and $A[\text{bias}, \text{bias}] \approx 1000$. Naive injection ($b \leftarrow b + 0.5$) yields only 0.05\% change. Confidence scaling ensures the bias is mathematically significant:

\begin{align}
\theta[\text{bias}] &= \frac{b[\text{bias}]}{A[\text{bias}, \text{bias}]} \\
\text{To shift } \theta \text{ by } \Delta: \quad &b_{\text{new}}[\text{bias}] = b[\text{bias}] + A[\text{bias}, \text{bias}] \times \Delta
\end{align}

\paragraph{Example:}
\begin{itemize}
    \item \textbf{Mixtral-8x7B (Fast/Cheap)}: 
    \begin{itemize}
        \item $\Delta_{\text{speed}} = +0.5$
        \item Effect: Biases toward 94.2\% Easy Cluster
        \item Result: High selection probability for routine tasks
    \end{itemize}
    
    \item \textbf{GPT-4-turbo (Slow/Expensive)}:
    \begin{itemize}
        \item $\Delta_{\text{speed}} = -0.5$
        \item Effect: Discourages selection (reserve for hard tasks)
        \item Result: Selected only when utility significantly compensates for cost penalty
    \end{itemize}
\end{itemize}

\subsubsection{Integration with Cost-Aware UCB}

The bias injection integrates with the cost-aware selection criterion:

\begin{equation}
m_t = \argmax_{m \in \mathcal{M}} \left[ \underbrace{\boldsymbol{\theta}_m^\top \mathbf{x}_t}_{\text{quality}} + \underbrace{\alpha \sqrt{\mathbf{x}_t^\top \mathbf{A}_m^{-1} \mathbf{x}_t}}_{\text{exploration}} - \underbrace{\lambda \cdot c_m}_{\text{cost penalty}} \right]
\label{eq:cost_aware_ucb}
\end{equation}

where:
\begin{itemize}
    \item $\boldsymbol{\theta}_m^\top \mathbf{x}_t$: Expected reward (includes bias injection)
    \item $c_m$: Normalized cost of model $m$ 
    \item $\lambda \geq 0$: Cost sensitivity parameter
\end{itemize}

This ensures business logic (speed profiles) and economic constraints (cost penalties) are jointly optimized with learned quality preferences.

\section{System Implementation}

\subsection{Three-Layer Architecture}

The complete initialization system operates in three layers:

\begin{enumerate}
    \item \textbf{Layer 1: Static Warm-Start}
    \begin{itemize}
        \item Load $\mathbf{A}_m$, $\mathbf{b}_m$ from 80k RouteLLM battles
        \item Apply Bayesian regularization: $\mathbf{A}_m \leftarrow \mathbf{A}_m + \lambda \mathbf{I}$
        \item Occurs in \texttt{CostAwareLinUCBRouter.\_\_init\_\_()}
    \end{itemize}
    
    \item \textbf{Layer 2: Semantic Transfer}
    \begin{itemize}
        \item Find semantic neighbor via DNA matching (Eq.~\ref{eq:neighbor_selection})
        \item Extract preferences: $\boldsymbol{\theta}_{\text{neighbor}} = \mathbf{A}_{m^*}^{-1} \mathbf{b}_{m^*}$ (Eq.~\ref{eq:preference_extraction})
        \item Reset confidence: $\mathbf{A}_{\text{new}} = \lambda \mathbf{I}$ (Eq.~\ref{eq:A_reset})
        \item Scale prior: $\mathbf{b}_{\text{new}} = \lambda \cdot \boldsymbol{\theta}_{\text{neighbor}} \cdot n_{\text{effective}}$ (Eq.~\ref{eq:b_initialization})
        \item Occurs in \texttt{BanditRouter.admix\_theta\_from\_neighbors()}
    \end{itemize}
    
    \item \textbf{Layer 3: T-Shirt Sizing}
    \begin{itemize}
        \item Apply speed biases: $b_m[\text{bias}] \leftarrow b_m[\text{bias}] + A_m[\text{bias}, \text{bias}] \times \Delta_{\text{speed}}$ (Eq.~\ref{eq:bias_injection})
        \item Occurs in \texttt{BanditRouter.create()}
    \end{itemize}
\end{enumerate}

\subsection{Algorithm: Complete Initialization Protocol}

\begin{algorithm}[t]
\caption{Three-Layer Initialization Protocol}
\label{alg:initialization}
\begin{algorithmic}[1]
\STATE \textbf{Input:} Models $\mathcal{M}$, warmup priors, registry, encoder
\STATE \textbf{Output:} Initialized router with warm-start
\STATE
\STATE \texttt{// Layer 1: Static Warm-Start}
\FOR{$m \in \mathcal{M}$}
    \STATE $\mathbf{A}_m \leftarrow \texttt{warmup\_priors}['A'][m]$
    \STATE $\mathbf{b}_m \leftarrow \texttt{warmup\_priors}['b'][m]$
    \STATE $\mathbf{A}_m \leftarrow \mathbf{A}_m + \lambda \mathbf{I}$ \COMMENT{Regularization}
\ENDFOR
\STATE
\STATE \texttt{// Layer 3: T-Shirt Sizing}
\FOR{$m \in \mathcal{M}$}
    \STATE $\text{speed} \leftarrow \texttt{registry}[m][\text{``speed\_profile''}]$
    \STATE $\Delta \leftarrow \begin{cases} +0.5 & \text{if speed = ``fast''} \\ -0.5 & \text{if speed = ``slow''} \\ 0.0 & \text{otherwise} \end{cases}$
    \STATE $\text{conf} \leftarrow \mathbf{A}_m[\text{bias}, \text{bias}]$
    \STATE $b_m[\text{bias}] \leftarrow b_m[\text{bias}] + \text{conf} \times \Delta$
\ENDFOR
\STATE
\STATE \textbf{return} Initialized router
\STATE
\STATE \texttt{// Layer 2: Semantic Transfer (on-demand)}
\FUNCTION{AddModel}{$m_{\text{new}}$, capabilities, speed}
    \STATE $\text{dna} \leftarrow \texttt{get\_dna}(m_{\text{new}}, \text{capabilities}, \text{speed})$
    \STATE $m^*, \text{sim} \leftarrow \texttt{find\_neighbor}(\text{dna}, \mathcal{M})$
    \IF{$\text{sim} > 0.5$}
        \STATE $\boldsymbol{\theta}_{\text{neighbor}} \leftarrow \mathbf{A}_{m^*}^{-1} \mathbf{b}_{m^*}$
        \STATE $n_{\text{eff}} \leftarrow \begin{cases} 5.0 & \text{if sim} > 0.8 \\ 3.0 & \text{if sim} > 0.6 \\ 1.0 & \text{otherwise} \end{cases}$
        \STATE $\mathbf{A}_{\text{new}} \leftarrow \lambda \mathbf{I}$ \COMMENT{Reset confidence}
        \STATE $\mathbf{b}_{\text{new}} \leftarrow \lambda \cdot \boldsymbol{\theta}_{\text{neighbor}} \cdot n_{\text{eff}}$ \COMMENT{Transfer preferences}
    \ELSE
        \STATE $\mathbf{A}_{\text{new}} \leftarrow \lambda \mathbf{I}$, $\mathbf{b}_{\text{new}} \leftarrow \mathbf{0}$ \COMMENT{Cold start}
    \ENDIF
    \STATE Apply Layer 3 (T-shirt sizing) to $m_{\text{new}}$
    \STATE $\mathcal{M} \leftarrow \mathcal{M} \cup \{m_{\text{new}}\}$
\ENDFUNCTION
\end{algorithmic}
\end{algorithm}

\section{Empirical Validation}

\subsection{Experimental Setup}

\paragraph{Dataset:} N=1,871 real prompts from production workload
\begin{itemize}
    \item Training: 1,121 prompts (dev set)
    \item Evaluation: 750 prompts (holdout set)
    \item Models: GPT-4-turbo (expensive) + Mixtral-8x7B (cheap)
\end{itemize}

\paragraph{Domain Shift:} Population Stability Index (PSI) = 0.42 between RouteLLM training data and production workload, indicating significant distribution shift.

\paragraph{Baselines:}
\begin{itemize}
    \item \textbf{Static GPT-4}: Always route to GPT-4-turbo
    \item \textbf{Static Mixtral}: Always route to Mixtral-8x7B
    \item \textbf{RouteLLM-MF}: Matrix factorization router~\cite{ong2024routellm}
    \item \textbf{BanditGPT (Cold-Start)}: LinUCB with $\mathbf{A}_m = \lambda \mathbf{I}$
    \item \textbf{BanditGPT (Warm-Start)}: Our three-layer initialization
\end{itemize}

\subsection{Results}

\begin{table}[t]
\centering
\caption{Performance at Production Quality (Reward = 0.90)}
\label{tab:main_results}
\begin{tabular}{@{}lccc@{}}
\toprule
Method & Avg Cost & Cost Reduction & Convergence \\
\midrule
Static GPT-4 & \$0.005500 & --- & Instant \\
Static Mixtral & \$0.000375 & 93.2\% & Instant \\
RouteLLM-MF & \$0.001834 & 66.7\% & N/A \\
\midrule
BanditGPT (Cold) & \$0.001250 & 77.3\% & 850 samples \\
\textbf{BanditGPT (Warm)} & \textbf{\$0.000423} & \textbf{92.3\%} & \textbf{200 samples} \\
\bottomrule
\end{tabular}
\end{table}

\paragraph{Key Findings:}
\begin{enumerate}
    \item \textbf{Cost Efficiency}: Warm-start achieves 92.3\% cost reduction, matching Static Mixtral's cost while maintaining 0.90 quality (vs. Mixtral's 0.78 quality).
    
    \item \textbf{Convergence Speed}: Warm-start converges in 200 samples vs. 850 for cold-start (4.25$\times$ faster).
    
    \item \textbf{Domain Shift Recovery}: Despite PSI=0.42, warm-start successfully adapts through multi-tiered initialization.
    
    \item \textbf{Cluster Exploitation}: Achieves 94.2\% routing to Mixtral for Easy Cluster while reserving GPT-4 for 5.8\% Hard Cluster.
\end{enumerate}

\subsection{Ablation Study}

\begin{table}[t]
\centering
\caption{Ablation: Impact of Each Layer}
\label{tab:ablation}
\begin{tabular}{@{}lccc@{}}
\toprule
Configuration & Avg Cost & Avg Reward & Convergence \\
\midrule
No warmup (cold) & \$0.001250 & 0.89 & 850 samples \\
Layer 1 only & \$0.000650 & 0.90 & 400 samples \\
Layer 1 + 3 & \$0.000480 & 0.90 & 250 samples \\
\textbf{All layers (1+2+3)} & \textbf{\$0.000423} & \textbf{0.90} & \textbf{200 samples} \\
\bottomrule
\end{tabular}
\end{table}

\paragraph{Analysis:}
\begin{itemize}
    \item \textbf{Layer 1 (80k battles)}: Provides 48\% cost improvement vs. cold-start
    \item \textbf{Layer 3 (T-shirt sizing)}: Additional 26\% improvement through speed biases
    \item \textbf{Layer 2 (Semantic transfer)}: Enables dynamic model addition without retraining
\end{itemize}

\subsection{Semantic Transfer Validation}

We evaluate semantic transfer by dynamically adding ``anthropic/claude-3.5-haiku'' after initial training:

\begin{table}[h]
\centering
\caption{Semantic Transfer: Haiku Addition}
\label{tab:transfer}
\begin{tabular}{@{}lcccc@{}}
\toprule
Method & Neighbor & Similarity & Samples to 90\% & Final Cost \\
\midrule
Cold-start & --- & --- & 650 & \$0.000520 \\
Naive transfer & Sonnet & 0.92 & 120 & \$0.000510 \\
$\boldsymbol{\theta}$-transfer & Sonnet & 0.92 & 85 & \$0.000390 \\
\bottomrule
\end{tabular}
\end{table}

\paragraph{Observations:}
\begin{itemize}
    \item Naive transfer (copy $\mathbf{A}$ and $\mathbf{b}$) requires 120 samples due to tight confidence intervals
    \item $\boldsymbol{\theta}$-transfer (reset $\mathbf{A}$) requires only 85 samples and achieves 23\% lower final cost
    \item This validates the ``confident transfer trap'' hypothesis
\end{itemize}

\section{Related Work}

\subsection{Contextual Bandits}

LinUCB~\cite{li2010contextual} provides theoretical guarantees for linear reward models but assumes cold-start initialization. Thompson Sampling~\cite{agrawal2013thompson} offers Bayesian exploration but requires expensive posterior sampling.

\subsection{Transfer Learning in Bandits}

Prior work on bandit transfer~\cite{lazaric2012transfer, bouneffouf2019survey} typically assumes similar source and target distributions. Our work explicitly addresses the ``confident transfer trap'' through $\boldsymbol{\theta}$-only transfer.

\subsection{LLM Routing}

RouteLLM~\cite{ong2024routellm} trains offline routers using matrix factorization but cannot adapt online. FrugalGPT~\cite{chen2023frugalgpt} uses cascade strategies but lacks contextual adaptation. Our work combines offline learning (warmup) with online adaptation (semantic transfer).

\section{Discussion}

\subsection{Theoretical Guarantees}

The $\boldsymbol{\theta}$-transfer protocol preserves LinUCB's regret guarantees. Since we reset $\mathbf{A}_{\text{new}} = \lambda \mathbf{I}$, the effective sample count is bounded:

\begin{equation}
N_{\text{effective}} = n_{\text{effective}} \ll N_{\text{neighbor}}
\end{equation}

This ensures the new model has sufficient exploration potential to achieve logarithmic regret:

\begin{equation}
R(T) = O\left( d\sqrt{T \log T} \right)
\end{equation}

\subsection{Production Considerations}

\paragraph{Computational Overhead:} Semantic DNA matching requires one forward pass through a sentence transformer (≈10ms). This is negligible compared to LLM inference latency (≈1000ms).

\paragraph{Memory Footprint:} Each model requires $O(d^2)$ memory for $\mathbf{A}_m$ (d=24 → 2.3 KB per model). For 100 models, total overhead is 230 KB.

\paragraph{Scalability:} The three-layer architecture scales linearly with the number of models. Warm-start loading is one-time (initialization), semantic transfer is on-demand (new model registration), and T-shirt sizing is $O(K)$ where $K$ is the model count.

\section{Conclusion}

We presented a production-grade contextual bandit router that achieves 92\% cost reduction at production quality through a three-layer initialization architecture:

\begin{enumerate}
    \item \textbf{Static Warm-Start}: Load 80k battle priors for immediate expertise
    \item \textbf{Semantic Transfer}: $\boldsymbol{\theta}$-only transfer avoids confident transfer trap
    \item \textbf{T-Shirt Sizing}: Confidence-scaled bias injection for business logic
\end{enumerate}

Our approach demonstrates robust recovery from domain shift (PSI=0.42) and enables dynamic model addition without retraining. The mathematical grounding ensures theoretical guarantees while the empirical validation on 1,871 real prompts confirms production viability.

\paragraph{Future Work:} Extensions include multi-objective optimization (quality + cost + latency), hierarchical clustering for finer-grained routing, and meta-learning for automatic hyperparameter tuning ($n_{\text{effective}}$, $\Delta_{\text{speed}}$).

\bibliographystyle{ACM-Reference-Format}
\begin{thebibliography}{9}

\bibitem{li2010contextual}
Lihong Li, Wei Chu, John Langford, and Robert E. Schapire.
\newblock A contextual-bandit approach to personalized news article recommendation.
\newblock In \emph{WWW}, pages 661--670, 2010.

\bibitem{agrawal2013thompson}
Shipra Agrawal and Navin Goyal.
\newblock Thompson sampling for contextual bandits with linear payoffs.
\newblock In \emph{ICML}, pages 127--135, 2013.

\bibitem{lazaric2012transfer}
Alessandro Lazaric.
\newblock Transfer in reinforcement learning: A framework and a survey.
\newblock In \emph{Reinforcement Learning}, pages 143--173. Springer, 2012.

\bibitem{bouneffouf2019survey}
Djallel Bouneffouf and Irina Rish.
\newblock A survey on practical applications of multi-armed and contextual bandits.
\newblock \emph{arXiv preprint arXiv:1904.10040}, 2019.

\bibitem{ong2024routellm}
Isaac Ong et al.
\newblock RouteLLM: Learning to route LLMs with preference data.
\newblock \emph{arXiv preprint arXiv:2406.18665}, 2024.

\bibitem{chen2023frugalgpt}
Lingjiao Chen et al.
\newblock FrugalGPT: How to use large language models while reducing cost and improving performance.
\newblock \emph{arXiv preprint arXiv:2305.05176}, 2023.

\bibitem{reimers2019sentence}
Nils Reimers and Iryna Gurevych.
\newblock Sentence-BERT: Sentence embeddings using Siamese BERT-networks.
\newblock In \emph{EMNLP}, 2019.

\end{thebibliography}

\end{document}

